Le lien électrofaible part du même opérateur angulaire qui a produit les constitutives du vide, mais utilise une autre fonctionnelle.\par

Pour obtenir \(\mu\), \(\varepsilon\), \(Z\) et \(c\), il fallait conserver le spectre complet. Il fallait distinguer les deux feuillets.\par

Le déterminant fait l’inverse :\par

\[
D_A=\det T=q_+q_-.
\]

En multipliant les valeurs propres, le déterminant intègre les deux orientations dans une grandeur angulaire commune et invariante. Il publie la partie paire et réserve l’identité du feuillet dans la mémoire orientée.\par

C’est pourquoi \(D_A\) est pair sous l’inversion\par

\[
C^*\longmapsto-C^*.
\]

Le secteur \(W/Z\) possède précisément cette structure. Les routes de \(W\) et de \(Z\) sont adjacentes dans la coordonnée angulaire principale et partagent la même coordonnée orientée. Lorsque l’on forme le quotient, la contribution de \(C^*\) s’annule.\par

Il reste\par

\[
\left(\frac{m_W}{m_Z}\right)^2
=
D_A,
\]

et donc\par

\[
\sin^2\theta_W
=
1-D_A.
\]

L’identité révèle que l’angle de Weinberg est le défaut complémentaire du déterminant angulaire.\par

Le déterminant mesure la partie commune qui subsiste lorsque les deux feuillets sont rabattus sur leur produit. \(1-D_A\) mesure la partie complémentaire nécessaire au mélange.\par

Nous pourrions le dire ainsi : le vide conserve toute la réponse ; le secteur électrofaible utilise la compression paire de cette réponse.\par

Le même antécédent angulaire produit deux fonctions différentes :\par

\begin{itemize}[leftmargin=*,itemsep=0.35em]
\item le calcul fonctionnel complet engendre les relations constitutives ;
\item le déterminant engendre l’invariant de mélange.
\end{itemize}

Cela explique pourquoi le vide et l’angle de Weinberg sont deux réalisations fonctionnelles distinctes d’un même antécédent angulaire.\par

À partir de là, les couplages \(g\) et \(g'\) sont, au niveau déclaré, les deux projections du même couplage électromagnétique rationalisé sur les directions sinus et cosinus du mélange :\par

\[
g^2
=
\frac{4\pi\alpha_{\rm HMT}}{1-D_A},
\qquad
g'^2
=
\frac{4\pi\alpha_{\rm HMT}}{D_A}.
\]

L’identité custodiale à l’arbre\par

\[
\rho_{\rm EW}^{(0)}=1
\]

exprime que la même partition angulaire organise de manière cohérente masse et couplage.\par

La constante de Fermi apparaît ensuite comme la réponse effective à basse énergie du canal chargé :\par

\[
G_F^{(0)}
=
\frac{\pi\alpha_{\rm HMT}}
{\sqrt2(1-D_A)m_W^2}.
\]

La constante de Fermi apparaît comme une composition de la structure fine, du défaut angulaire et de l’échelle du porteur \(W\).\par

Le secteur de Higgs introduit une différence décisive. Sa route possède une coordonnée orientée différente de celle de la route de \(W\), de sorte que \(C^*\) persiste. L’autocouplage scalaire publie l’orientation que le déterminant réserve.\par

Le secteur de jauge est pair : il utilise \(D_A\).\par
Le secteur scalaire est orienté : il retient \(C^*\).\par

Cette division est physiquement suggestive. Les champs de jauge organisent la partie commune de la transformation ; le secteur scalaire conserve la mémoire de l’orientation brisée.\par

Les corrections radiatives appartiennent à un niveau ultérieur et doivent être construites prospectivement à partir des boucles, états et échelles internes, avec \(G_F\) comme sortie et évaluation terminale. Le squelette causal est déjà présent : vide, mélange, couplages et réponse de Fermi procèdent de différentes lectures du même état angulaire enrichi.\par

